% !TeX root = ../papers/nc-p3.tex
\section{四中心切断を持つ窓に対する次数16の完全判定}
\label{sec:central-window-cutoff}
この判定は成分、反対階数、PBW性を仮定しない。四窓に六つずつの二次関係を与え、
有限双対の標準Frobenius検査を先に済ませる。
以下の次数は全て元の経路長である。
\begin{theorem}\label{thm:central-window-cutoff}
折り畳み環$B$に四つの中心四次式$x_1,\ldots,x_4$があり、
$Q=B/(x_1,\ldots,x_4)$が有限次元とする。このとき標準Koszul左右AS正則性は
次の有限条件と同値である。
\begin{enumerate}
\item 各始点の$Q$の次数$0$から13までの次元は$(1+t+t^2+t^3)^4$の係数に一致する。
\item 各始点の$B$の次数$n\le16$の次元は$\binom{n+3}{3}$であり、
左右の標準線型複体は内部次数16まで単純加群を解消する。
\end{enumerate}
特に$Q_{13}=0$だけで有限性の仮定を認証でき、以後の判定は経路長16を超えない。
適切な切断を持つ標準AS窓の近くで、四切断を含む有限Frobenius候補の
パラメータ空間上、これらの認証条件は開く。
\end{theorem}
\begin{proof}
$P=k[X_1,\dots,X_4]$, $\deg X_i=4$を$x_i$で作用させる。
$Q$が有限なら次数帰納で$B$は$P$上有限である。まず$B$を標準ASとする。
中央Koszul複体$K(X;B)$の正のホモロジーも含め全ホモロジーは$P_+$で消え、
有限次元の次数付き$B$加群である。有限次元次数付き加群は頂点単純の移動による
有限列を持つ。AS双対性と長完全列により、非零なそのような加群$H$について
$\operatorname{Ext}_B^p(H,B)$は$p=4$だけで非零となる。
各項が$B$自由な中央Koszul複体に$\operatorname{Hom}_B(-,B)$を施すと、
スペクトル系列の$\operatorname{Ext}^p_B(H_q,B)$は全次数$p+q$へ収束する。
$p=4$の一列しかなく、複体の長さは4なので、$q>0$の$H_q$は零である。
従って$B$の$P$上の高次Torは消え、極小次数付き自由分解から$B$は$P$自由である。
各始点の標準Hilbert級数に$(1-t^4)^4$を掛ければ第一条件を得る。第二条件は正則性から従う。

逆に二条件を仮定する。$Q$の各始点・終点の斉次基底を経路から選び、
その語$w$に対応する記号を基底とした自由$P$加群$F$を取る。最長は12である。
$Q_{13}=0$と一次生成性より全高次も零であり、評価$F\to B$は全射である。
第一条件から$F$の各始点級数は$(1-t)^{-4}$で、第二条件によりこの評価は
次数16まで同型となる。ここから全次数の同型が従うことを示す。
各矢$a$の左作用を、基底語への積$aw$の一意な座標表示で定義し、$P$線型に
$L_a:F\to F$へ延ばす。端点が合わない積は零とする。この操作は長さ13までで済む。
評価と矢は可換である。関係$r$の作用$r(L)$は各基底上で評価が零で、
長さは高々14なので$F$でも零である。全ての中心切断についても
$x_i(L)-X_i$は各基底上で長さが高々16、評価が零なので作用は零である。
従って$F$は$B$加群である。頂点単位の和に語を作用させれば、
対応する基底記号は次数12までの単射性により全て回収され、
さらに$X_i$の作用は$x_i$で回収される。これで$B\to F$が評価の逆写像となる。
特に標準Hilbert関数と自由性が全次数で確定した。

左右どちらかの頂点単純に対する線型複体を、増大写像付きで$C_\bullet$と書く。
各項は$P$自由で、係数空間の次数移動は0から4なので自由生成元の次数は高々16である。
同じ単純の$P$自由分解は四変数の通常のKoszul複体$D_\bullet$であり、
生成元次数は$0,4,8,12,16$である。増大写像の恒等を持ち上げて
$D_\bullet\to C_\bullet$を作れる。実際、各自由生成元の像を選ぶ際の閉元は
内部次数16以下にあり、仮定した範囲の完全性で境界として解ける。
比較写像の錐は次数16まで完全で、自由生成元に16を超えるものがない。
定数の単元成分を逐次縮約して極小化する。非零な項が残れば、内部次数が最小の
自由生成元は正次数係数の微分で零に移り、境界にもなれない。
この生成元の次数は16以下であり矛盾する。従って錐は可縮で、もとの複体が
全次数で単純を解消する。有限Frobenius性と合わせて求める標準AS性が従う。

最後に正則な基点で、次数16までの左右線型複体の単元小行列を開く。
次数別に帰納すれば、前次数の商空間が束であり、複体左端からの分裂完全性で
次の次数の商も標準階数の束となる。同じ範囲の中央Koszul複体も基点で
正の位置が完全であるから分裂する小行列を選べる。従って次数13までのその
余核の次元も所要値となる。有限個の開条件だけで定理を適用できる。
\end{proof}

\path{computations/verify_central_window_cutoff.py}の\texttt{certify}は有理係数で
この判定を実施する。四つの候補切断の中心性を長さ5で確認し、六関係と
四中心切断を挿入した語空間を長さ13まで構成する。$Q_{13}\ne0$で有限性が
別途未認証なら切断未確認を返す。切断を確認できた場合には、語の逐次商と
双対交差空間の一文字縮約から左右の複体の像階数を長さ16まで比較し、
正則・非正則のどちらも有限に判定する。
四切断を存在量化する場合、中心性とこれらの有限行列条件は可構成で、
標数零の代数閉体上でパラメータを消去できる。全標準AS点がこのような四切断を
持つことは必要とせず、その被覆は別の問題である。

\begin{corollary}
中心四次式を指定しない入力でも、$B$が四次中心全体の生成する部分環上有限である
という範囲では、同じ次数16まででAS正則性を決定できる。
\end{corollary}
\begin{proof}
長さ5の可換方程式で四次中心の基底を求め、独立な不定係数で四本の一般線型結合を取る。
標準ASなら$B$の成長次元は4で、有限性によりその中心部分環のKrull次元も4である。
この可換環は換算次数1で生成されるため、無限体上で一般の四線型式が
斉次Noether正規化を与える。従ってこの四本の剰余は有限次元であり、
定理の必要条件から長さ13で必ず零になる。よって有理関数体上でその次数が
非零なら、前記の有限性の仮定の下で非ASである。零なら定理の二条件を適用し、
体の忠実平坦な拡大に沿う完全性の降下で結論を得る。
この有限性には、全四次中心で割った商が長さ13で零という直接確認できる十分条件もある。
\end{proof}
